\par\noindent\textbf{(1)} 良导体内∇×H≈σE，∇×E=iωμH（e\textasciicircum{}\{-iωt\}约定）。联立有∇²E=−iωμσE。设E\_x(z)=E0 e\textasciicircum{}\{−γz\}，得γ²=−iωμσ，取Reγ>0根γ=(1−i)/δ，δ=√[2/(μσω)]，所以场幅向内以e\textasciicircum{}\{−z/δ\}衰减。
\par\noindent\textbf{(2)} 对E\_x(z)=E0 exp[−γz]，其旋度的y分量是∂\_zE\_x=−γE\_x。频域Faraday方程∇×E=iωμH给H\_y=(1/(iωμ))∂\_zE\_x=−γE\_x/(iωμ)=iγE\_x/(ωμ)，其中γ=(1−i)/δ。于是E×H沿+z输送正能流；场幅e\textasciicircum{}\{−z/δ\}衰减而相位随深度旋转，表面E与H相差45°。
\par\noindent\textbf{(3)} 单位表面平均Joule功P\_A=∫\_0∞σ|E(z)|²/2 dz=σδ|E0|²/4。由表面阻抗Z\_s=E\_t/H\_t=(1−i)√[ωμ/(2σ)]（取方向一致的复比）得表面电阻R\_s=ReZ\_s=1/(σδ)=√(ωμ/(2σ))，并有P\_A=R\_s|H\_t|²/2，证明输入功转为体热。
\par\noindent\textbf{(4)} ω增大使δ∝ω\textasciicircum{}\{−1/2\}减小、R\_s∝√ω增大，电流越集中导体表面、交流损耗越强。σ→∞时δ→0、R\_s→0，理想导体切向电场消失；在纳米尺度平均自由程可与δ相比，局域欧姆定律失效。
